\section{Additional Scalability Measurements}
\label{app:b3-scalability}
\label{app:b3-concurrency}

\noindent\textbf{Protocol scaling.}
Table~\ref{tab:b3-protocol-scaling} reports the ten CPU protocol configurations
underlying the scalability curves: five call counts and two lifecycle-event
densities. Construction includes in-memory protocol generation and identity
initialization; audit checks the resulting signature-profile transcript.
Model inference, evidence persistence, ledger I/O, and proof generation are excluded.
Each configuration is measured once. The shaded envelopes in the main figure
therefore span configurations, rather than repeated-run uncertainty.
At 512 calls and 4,096 events, construction and audit take 764.82 and 864.69\,ms,
respectively.
\begin{table*}[!t]
\centering
\caption{Protocol construction and audit costs (one measurement per configuration).}
\label{tab:b3-protocol-scaling}
\footnotesize
\renewcommand{\arraystretch}{1.10}
\begin{tabular*}{\textwidth}{@{\hspace{4pt}}c@{\extracolsep{\fill}}cccc@{\hspace{4pt}}}
\toprule
\textbf{Calls/request} & \multicolumn{2}{c}{\textbf{4 events/call}} & \multicolumn{2}{c}{\textbf{8 events/call}} \\
\cmidrule(lr){2-3}\cmidrule(lr){4-5}
 & \textbf{Build (ms)} & \textbf{Audit (ms)} & \textbf{Build (ms)} & \textbf{Audit (ms)} \\
\midrule
1 & 3.48 & 1.82 & 2.16 & 2.22 \\
8 & 9.52 & 10.18 & 12.59 & 13.76 \\
32 & 35.94 & 38.87 & 74.94 & 71.81 \\
128 & 160.79 & 154.50 & 190.25 & 212.52 \\
512 & 566.98 & 620.65 & 764.82 & 864.69 \\
\bottomrule
\end{tabular*}
\end{table*}


\noindent\textbf{Serving concurrency.}
Table~\ref{tab:b3-concurrency} reports the original Qwen3-1.7B concurrency
measurements. Each scheme uses one measured round per configuration, with 16
requests at concurrency 1 and 128 at concurrency 8 and 32; warmups are excluded.
The concurrency-8 entries are the original round, rather than the later
three-round averages; these observations are reused in the cross-model cohort.
Each request comprises eight sequential calls, with
2,048 input tokens per call and an output limit of 256 tokens.
Latency is mean request processing time, excluding waiting for a worker slot.
Request and output throughput
are completed requests and generated tokens, respectively, divided by batch wall time.
Online non-LLM time is $T_{\mathrm{online}}-T_{\mathrm{llm}}$, where
$T_{\mathrm{llm}}$ accumulates model-service HTTP wall time.
Actual output length is included because the realized generation workload differs
across concurrency settings; request throughput alone does not isolate concurrency.
% Requires booktabs and xcolor with the table option (or colortbl).
\begin{table*}[!t]
\centering
\caption{Concurrency measurements for Qwen3-1.7B (one original round per configuration).}
\label{tab:b3-concurrency}
\footnotesize
\renewcommand{\arraystretch}{1.10}
\begin{tabular*}{\textwidth}{@{\hspace{4pt}}l@{\extracolsep{\fill}}cccccc@{\hspace{4pt}}}
\toprule
\textbf{Scheme} & \textbf{Requests} & \shortstack{\textbf{Latency}\\\textbf{(s)}} & \shortstack{\textbf{Throughput}\\\textbf{(req/s)}} & \shortstack{\textbf{Output rate}\\\textbf{(tokens/s)}} & \shortstack{\textbf{Online non-LLM}\\\textbf{(ms)}} & \shortstack{\textbf{Output length}\\\textbf{(tokens/call)}} \\
\midrule
\rowcolor{gray!12}
\multicolumn{7}{l}{\textbf{Concurrency 1}} \\
Vanilla & 16 & 6.893 & 0.145 & 65.51 & 1.18 & 56.45 \\
SignedLog & 16 & 6.995 & 0.143 & 64.56 & 3.57 & 56.45 \\
SignedLog+ & 16 & 6.999 & 0.143 & 64.51 & 14.03 & 56.45 \\
\textbf{\textsc{TokenProv}} & 16 & 7.033 & 0.142 & 64.20 & 13.33 & 56.45 \\
\midrule
\rowcolor{gray!12}
\multicolumn{7}{l}{\textbf{Concurrency 8}} \\
Vanilla & 128 & 11.365 & 0.703 & 192.49 & 1.30 & 34.24 \\
SignedLog & 128 & 14.719 & 0.543 & 147.35 & 10.80 & 33.95 \\
SignedLog+ & 128 & 14.997 & 0.529 & 144.64 & 58.24 & 34.16 \\
\textbf{\textsc{TokenProv}} & 128 & 11.409 & 0.700 & 191.68 & 52.46 & 34.23 \\
\midrule
\rowcolor{gray!12}
\multicolumn{7}{l}{\textbf{Concurrency 32}} \\
Vanilla & 128 & 67.262 & 0.472 & 245.10 & 1.00 & 64.91 \\
SignedLog & 128 & 89.193 & 0.356 & 184.74 & 11.17 & 64.91 \\
SignedLog+ & 128 & 72.323 & 0.438 & 227.59 & 86.04 & 64.91 \\
\textbf{\textsc{TokenProv}} & 128 & 59.350 & 0.535 & 277.84 & 85.01 & 64.90 \\
\bottomrule
\end{tabular*}
\end{table*}

